\section{Conservative Generator Correspondences}\label{app:conservativecalculations}
\status{explicit standard calculations and bounded reference models; none is a native Thalean gyro or Maxwell construction}
This appendix supplies the calculations behind Sections~37A--37D and the exact identities tested by the new verification script. The word ``proved'' identifies a deduction from declared hypotheses, not a novelty claim for standard linear algebra, frame geometry, or classical field theory.

\subsection{Finite isometries, moving forms, and the Cayley step}
For fixed $\mathsf G$, differentiation of the quadratic form gives (CD2). For a time-dependent form, applying the product rule to $X^\dagger\mathsf G(t)X$ yields (CD5); omitting $\dot{\mathsf G}$ is generally wrong. For example, $X(t)=e^tX(0)$ and $\mathsf G(t)=e^{-2t}I$ preserve the moving readout even though $\mathsf A=I$ is not skew relative to a fixed positive metric. The interpretation of a changing metric must itself be declared.

For real step $h$, write $B_\pm=I\pm h\mathsf A/2$. Under (CD2), expansion gives
\[
B_+^\dagger\mathsf G B_+
=\mathsf G+\frac{h^2}{4}\mathsf A^\dagger\mathsf G\mathsf A
=B_-^\dagger\mathsf G B_-.
\]
Since $B_+$ and $B_-$ commute, $U_h=B_-^{-1}B_+=B_+B_-^{-1}$, and multiplication by $B_-^{-\dagger}$ and $B_-^{-1}$ proves $U_h^\dagger\mathsf G U_h=\mathsf G$. For $\mathsf G\succ0$ a skew generator is similar to a skew-Hermitian matrix, so $B_-$ is invertible for real $h$. Also $U_h=I+h\mathsf A+O(h^2)$. Conservation of the form does not make this rational update identical to the exponential; the two have different finite-step phase advances.

The distinction between state recovery and form preservation survives exact arithmetic. The invertible map $2I$ changes every nonzero positive norm; the map $x\mapsto\|x\|e_1$ preserves that norm and identifies, for example, $e_1$ and $e_2$. These are counterexamples to opposite false implications. The identity is a separate negative control for claims of informative effect.

\subsection{The local rotational algebra and its scope}
For $\boldsymbol\Omega=(\Omega_x,\Omega_y,\Omega_z)$ define
\[
[\boldsymbol\Omega]_\times=
\begin{pmatrix}
0&-\Omega_z&\Omega_y\\
\Omega_z&0&-\Omega_x\\
-\Omega_y&\Omega_x&0
\end{pmatrix}.
\]
Every real three-dimensional skew matrix has exactly this form. If $Q(t)\in SO(3)$, differentiating $QQ^{\mathsf T}=I$ gives $\dot Q Q^{\mathsf T}=[\boldsymbol\omega_Q]_\times$. Hence $\mathbf n=Qe_3$ satisfies $\dot{\mathbf n}=\boldsymbol\omega_Q\times\mathbf n$. Conversely, specifying only $\mathbf n(t)$ leaves the axial component undetermined.

Tangency is weaker than a global rotational action. Fix $\mathbf a=e_3$ and take $F(\mathbf n)=(1+n_z)(e_3\times\mathbf n)$. Its values at $e_1$ and $e_2$ would force any constant generator to be $e_3$. At $\mathbf n=(3/5,0,4/5)$ that generator gives $(0,3/5,0)$, whereas $F(\mathbf n)=(0,27/25,0)$. Thus no constant matrix of this form acts on all these orientations. The law remains tangent and is covariant when $\mathbf a$ is rotated with $\mathbf n$.

For a constant generator, $Q(t)=\exp(t[\boldsymbol\Omega]_\times)$ preserves both $\|\mathbf n\|$ and $\boldsymbol\Omega\cdot\mathbf n$. This proves the conical orbit condition. A state-dependent or time-dependent generator still has tangent velocity but need not preserve any fixed cone.

\subsection{Stabilizer geometry and local phase}
Let $h_\theta=R_z(\theta)$ and $s=\operatorname{diag}(1,-1,-1)$. Then $\det s=1$, $se_3=-e_3$, and
\[
s^2=I,\qquad sh_\theta s=h_{-\theta}.
\]
The matrices fixing $e_3$ are the $h_\theta$. The matrices preserving the unoriented line $\R e_3$ form the two components $\{h_\theta\}\cup\{s h_\theta\}$, an embedded $O(2)$ in $SO(3)$. This proves the distinction between the oriented sphere $S^2$ and the unoriented base $\mathbb{RP}^2$. Restricting $h_\theta$ to an order-$n$ cyclic subgroup gives an abstract dihedral group of order $2n$. The native instances still require their own action maps.

For the local Euler chart used in (GY4), angular velocities add in the spatial frame as
\[
\boldsymbol\omega_Q
=\dot\varphi e_3+\dot\theta R_z(\varphi)e_2+\dot\phi\mathbf n.
\]
Taking the dot product with $\mathbf n=R_z(\varphi)R_y(\theta)e_3$ proves (GY4). A different transverse-frame convention changes the local fiber angle and its connection contribution together. There is no global smooth choice of a unit transverse vector at every point of $S^2$: it would be a nowhere-zero tangent vector field. This is why the frame bundle is not a global product, and why the local phase formula must not be treated as a globally defined independent scalar coordinate \citep{crainicbundles}.

The infinitesimal quotient in (GY2) is an isomorphism of vector spaces at a fixed axis. The global map $SO(3)\to S^2$ is not a group homomorphism. The axis stabilizer changes by conjugation with the axis; it is not an invisible normal subgroup of the full rotation action on every axis simultaneously.

\subsection{Torque and the regular heavy top}
Let $\mathbf L=\ell\mathbf n$ with $\ell>0$. Dotting $\dot{\mathbf L}=\boldsymbol\tau$ with $\mathbf n$ gives $\dot\ell=\mathbf n\cdot\boldsymbol\tau$. Subtracting that component gives $\ell\dot{\mathbf n}=\boldsymbol\tau_\perp$. The triple-product identity
\[
(\mathbf L\times\boldsymbol\tau)\times\mathbf L
=\ell^2\boldsymbol\tau-(\mathbf L\cdot\boldsymbol\tau)\mathbf L
\]
then proves (GY5). These identities hold about the stated fixed origin; replacing the origin or pivot with a moving one requires the appropriate additional balance terms.

For a symmetric heavy top with fixed pivot let $\mathbf a$ be its body symmetry axis, let $e_z$ point upward, and let $\mathbf a\cdot e_z=\cos\theta$ be constant in regular precession. Let $\Omega$ be its precession rate, $I_1$ its transverse moment about the pivot, and $L_3=I_3\omega_3$ its axial angular momentum. Kinematics and the inertia tensor give
\[
\boldsymbol\omega=\Omega(e_z-\cos\theta\,\mathbf a)+\omega_3\mathbf a,
\qquad
\mathbf L=I_1\Omega(e_z-\cos\theta\,\mathbf a)+L_3\mathbf a.
\]
For constant $\Omega,\theta,L_3$, use $\dot{\mathbf a}=\Omega e_z\times\mathbf a$ and the gravitational torque $Mgr\,e_z\times\mathbf a$. At nonzero tilt the exact balance becomes
\begin{equation}
\Omega(L_3-I_1\Omega\cos\theta)=Mgr,
\qquad
I_1\Omega^2\cos\theta-L_3\Omega+Mgr=0.
\tag{GY6}\label{eq:regulartop}
\end{equation}
Where the real slow branch exists and $|I_1\Omega\cos\theta|\ll|L_3|$, its leading value is $Mgr/L_3$. When $\cos\theta\ne0$ the quadratic's two roots have discriminant $L_3^2-4I_1Mgr\cos\theta$; negative discriminant excludes regular precession at that specified tilt. At $\cos\theta=0$ the equation reduces directly to $L_3\Omega=Mgr$. At exact vertical alignment the cross-product balance is degenerate and must not be obtained by cancelling $\sin\theta$. Initial conditions with nutation require the fuller top dynamics, not this regular-precession ansatz \citep{feynmanrotation}.

The mechanical input is the torque balance and inertia law. The computation explains its relation to the conservative orientation register but does not derive either input from a finite carrier.

\subsection{The periodic Maxwell operator and its domain}
Take $\mathbb T^3=(\R/2\pi\Z)^3$ and complexify the vector fields for Fourier analysis. Define the maximal periodic curl domain by
\[
\mathcal D(\mathscr C)
=\left\{u\in L^2(\mathbb T^3;\C^3):
\sum_{k\in\Z^3}|k\times\widehat u(k)|^2<\infty\right\}.
\]
On each Fourier mode, $\mathscr C(k)=i[k]_\times$ is Hermitian, since $[k]_\times^{\mathsf T}=-[k]_\times$. The direct-sum operator on this maximal domain is therefore self-adjoint. Consequently (EM3), on $\mathcal D(\mathscr C)\oplus\mathcal D(\mathscr C)$, is skew-adjoint and generates a unitary $L^2$ group. The real-field condition $\widehat\Psi(-k)=\overline{\widehat\Psi(k)}$ is preserved. This proves the precise periodic operator claim without assuming that all boundary realizations have the same adjoint.

The symbol obeys
\begin{equation}
\mathsf A_M(k)^2=c^2
\begin{pmatrix}
kk^{\mathsf T}-|k|^2I_3&0\\
0&kk^{\mathsf T}-|k|^2I_3
\end{pmatrix}.
\tag{EM7}\label{eq:maxwellsymbol}
\end{equation}
For $k\ne0$ the divergence-free condition removes the two longitudinal directions. The four-dimensional complex transverse sector has eigenvalues $+ic|k|$ and $-ic|k|$, each twice. The $k=0$ sector consists of spatially uniform static fields. None of these longitudinal or zero-frequency kernel directions is automatically a rotor phase fiber. The Fourier calculation establishes a conventional comparison spectrum, not a finite Thalean dispersion relation.

Integration by parts on a region with boundary instead gives, for smooth fields,
\[
\int_V u\cdot(\nabla\times v)\,\mathrm d^3x
-\int_V(\nabla\times u)\cdot v\,\mathrm d^3x
=-\int_{\partial V}\boldsymbol\nu\cdot(u\times v)\,\mathrm dA.
\]
Complex fields use $\overline u$ in the first slot. Boundary-domain conditions must remove the appropriate bilinear form for all admissible pairs, and maximality is needed for skew-adjointness rather than merely formal skew symmetry. Boundary realizations are part of the operator, not a footnote to it \citep{ellerkarabash2021}.

For real fields, multiply the two curl equations by $\epsilon_0\mathbf E$ and $\mathbf B/\mu_0$. The identity $\nabla\cdot(\mathbf E\times\mathbf B)=\mathbf B\cdot\nabla\times\mathbf E-\mathbf E\cdot\nabla\times\mathbf B$ gives (EM5). This explicit source and boundary accounting distinguishes global conservation from local redistribution.

\subsection{Reciprocity and duality are different operators}
Let $\mathsf D=\operatorname{diag}(\mathscr C,\mathscr C)$. Then $\mathsf A_M=c\mathsf J\mathsf D$, with $\mathsf J^\dagger=-\mathsf J$, $\mathsf J^2=-I$, and $[\mathsf J,\mathsf D]=0$. These facts prove that constant $e^{\vartheta\mathsf J}$ commutes with the evolution. They do not identify this internal rotation with spatial rotation or time translation.

The unsigned swap $\mathsf P=\left(\begin{smallmatrix}0&I\\I&0\end{smallmatrix}\right)$ instead obeys $\mathsf P\mathsf J\mathsf P=-\mathsf J$. Thus $\mathsf P\mathsf A_M\mathsf P=-\mathsf A_M$. For a conservative curl domain, changing the evolution's lower-left sign to positive gives a self-adjoint rather than skew-adjoint block. If curl has a nonzero eigenvalue $\lambda$, that altered block has real growth/decay eigenvalues $\pm c\lambda$. Norm preservation is lost. The relative sign is therefore part of the conservation law, not a decorative choice.

\subsection{A reference incidence complex, not a native construction}
For the finite cochain target of Appendix~\ref{app:horizon}, put $X=(e,b)$,
\[
\mathsf G_f=\begin{pmatrix}M_e&0\\0&M_b\end{pmatrix},\qquad
\mathsf A_f=\begin{pmatrix}
0&M_e^{-1}D_1^{\mathsf T}M_b\\-D_1&0
\end{pmatrix}.
\]
Then $\mathsf G_f\mathsf A_f=\left(\begin{smallmatrix}0&D_1^{\mathsf T}M_b\\-M_bD_1&0\end{smallmatrix}\right)$ is skew, so $\mathsf A_f^{\mathsf T}\mathsf G_f+\mathsf G_f\mathsf A_f=0$. The constraint operator
\[
\mathsf B_f=
\begin{pmatrix}D_0^{\mathsf T}M_e&0\\0&D_2\end{pmatrix}
\]
satisfies $\mathsf B_f\mathsf A_f=0$, using $D_1D_0=0$ and $D_2D_1=0$. With forcing $(-M_e^{-1}j,0)$, the electric constraint evolves by $-D_0^{\mathsf T}j$ and the magnetic constraint is unchanged; half the quadratic energy evolves by $-e^{\mathsf T}j$. Thus an independently supplied charge law $\dot\rho+D_0^{\mathsf T}j=0$ preserves electric Gauss compatibility. These are algebraic consequences of the supplied complex and pairings, not an electromagnetic interpretation of arbitrary cochains.

The included script constructs the oriented tetrahedral simplex complex directly from its vertices, edges, faces, and volume. It tests the two boundary-of-boundary identities, weighted skewness, constraint preservation, an exact Cayley step, and the source balance. It also tests a six-component Fourier symbol at a nonzero wavevector. Neither construction uses $G_{60}$, WXYZTI, or the native event-admission grammar. They are comparison harnesses for future native candidates, with their non-native status recorded in the generated report.

\subsection{Verification and limits of the tests}
\path{scripts/verify_conservative_correspondences.py} checks exact symbolic identities, rational examples, and negative controls. These include a tangent nonlinear non-isometry, a norm-preserving many-to-one map, an invertible nonconservative map, Euler norm drift, an unsigned reciprocal-sign failure, and a skew generator that violates a chosen divergence constraint. Exact finite tests complement the displayed general proofs; they do not prove continuum convergence, a universal information metric, source-selected dynamics, or a physical QR correspondence. The inherited verification scripts and all their native-status boundaries remain intact.
